\(\Q\cap[0,1]\) 在 Lebesgue 测度下的``长度''为 0——可数集可以零测。而 \([0,1]\) 本身的测度为 1。这说明零测集不等于空集，概率为 0 不等于不可能。这一区分必须从测度构造本身来理解。

前面的概率论把概率空间写成 \((\Omega,\mathcal F,P)\)，把随机变量写成可测函数，也已经使用单调收敛、控制收敛与条件期望。本单元补上这些概念背后的测度论结构。主线由六个问题组成：

\begin{enumerate}
\tightlist
\item
  \(\sigma\)-代数把可观察事件组织成封闭系统：从有限分割、Borel 集与随机变量生成的信息解释可测性。
\item
  测度把长度、计数与概率统一起来：用可数可加性推出集合极限，并区分零测集与不可能事件。
\item
  Lebesgue 积分先积简单函数再逼近：把面积、期望、无穷和与极限定理放进同一构造。
\item
  Radon--Nikodym 导数把``密度相对于谁''说清楚：统一概率密度、似然比与换测度，并处理没有普通密度的分布。
\item
  条件期望是对信息 \(\sigma\)-代数的投影：给出不依赖条件密度的定义，解释几乎处处唯一、塔式性质与最佳预测。
\item
  一致可积性阻止少量大值携带期望质量：说明何时能把依概率收敛升级为 \(L^1\) 收敛。
\end{enumerate}

初次接触概率时，先用离散求和与连续密度建立直觉是合理的；回到这里是为了让同一套语言能够覆盖混合分布、随机过程、重要性采样、martingale 和 Bellman 条件期望，而不再依赖``应该存在一个密度''的假定。

\subsection{一个具体算例}

设 \(\Omega=\{1,2,3,4,5,6\}\)（掷骰子），\(\sigma\)-代数 \(\mathcal F=2^\Omega\)，概率测度 \(P(A)=|A|/6\)。随机变量 \(X(\omega)=\omega\) 对 \(\mathcal F\) 可测。若改用较粗的 \(\sigma\)-代数 \(\mathcal G=\{\emptyset,\{1,3,5\},\{2,4,6\},\Omega\}\)（只区分奇偶），则 \(X\) 对 \(\mathcal G\) 不可测，但 \(Y(\omega)=\mathbf{1}_{\{\omega\text{为奇数}\}}\) 对 \(\mathcal G\) 可测。条件期望 \(E[X\mid\mathcal G]\) 在每个原子上是 \(X\) 的平均：对奇数原子 \(\{1,3,5\}\) 取 \((1+3+5)/3=3\)，对偶数原子 \(\{2,4,6\}\) 取 \((2+4+6)/3=4\)。

\subsection{怎么读这一章}

按``事件---大小---积分---换参考---加信息''推进。\(\sigma\)-代数规定哪些事件可谈，测度给它们大小，Lebesgue 积分累计函数值，Radon--Nikodym 导数比较两个测度，条件期望则只保留当前信息能看见的部分。最后用一致可积性控制期望不逃逸。

\subsection{本章篇目}

以下篇目按概念依赖排列；若从具体问题进入，也可以先打开最接近当前困惑的一篇，再沿篇内链接回补。

\begin{itemize}
\tightlist
\item \hyperref[sec:12-Real-Analysis-Support/05-Measure-Integration-and-Conditional-Expectation/01]{``\(\sigma\)-代数把可观察事件组织成封闭系统''}；
\item \hyperref[sec:12-Real-Analysis-Support/05-Measure-Integration-and-Conditional-Expectation/02]{``测度把长度、计数与概率统一起来''}；
\item \hyperref[sec:12-Real-Analysis-Support/05-Measure-Integration-and-Conditional-Expectation/03]{``Lebesgue 积分先积简单函数再逼近''}；
\item \hyperref[sec:12-Real-Analysis-Support/05-Measure-Integration-and-Conditional-Expectation/04]{``Radon--Nikodym 导数把``相对于谁''说清楚''}；
\item \hyperref[sec:12-Real-Analysis-Support/05-Measure-Integration-and-Conditional-Expectation/05]{``条件期望是对信息 \(\sigma\)-代数的投影''}；
\item \hyperref[sec:12-Real-Analysis-Support/05-Measure-Integration-and-Conditional-Expectation/06]{``一致可积性阻止期望质量逃逸''}。
\end{itemize}
